\section*{Bab \ref{ch:b1:wave-propagation} --- Perambatan Sinyal: Gelombang Berjalan dan Gelombang Berdiri}

\begin{solution}{exo:b1:wave-propagation:1}
$d = 340 \times 6.0 = \qty{2.0}{km}$. Bolak-balik: kedalamannya $= 1500 \times
0.40/2 = \qty{300}{m}$.
\end{solution}

\begin{solution}{exo:b1:wave-propagation:2}
$\lambda = 340/440 = \qty{0.77}{m}$; $\Delta\varphi = 2\pi \times
0.20/0.773 = \qty{1.6}{rad}$ ($93^\circ$); sefase untuk pisah sejauh
$n\lambda = n \times \qty{0.77}{m}$.
\end{solution}

\begin{solution}{exo:b1:wave-propagation:3}
$c = \sqrt{80/4.0\times10^{-3}} = \qty{141}{m/s}$; $\lambda = c/f =
\qty{0.71}{m}$.
\end{solution}

\begin{solution}{exo:b1:wave-propagation:4}
$444 - 440 = 4$ layangan tiap sekon; periode selubungnya $1/4 = \qty{0.25}{s}$.
Mengendurkannya menurunkan senarnya: \qty{441}{Hz} (mencapai \qty{439}{Hz}
berarti harus melewati kesenyapan pada \qty{440}{Hz}); periksalah dengan
mengendurkannya sedikit lagi: layangannya harus lenyap, lalu kembali.
\end{solution}

\begin{solution}{exo:b1:wave-propagation:5}
$f_1 = c/2L = 240/1.2 = \qty{200}{Hz}$; $400$, \qty{600}{Hz}. Moda ketiga:
simpul di $0$, $0.20$, $0.40$, \qty{0.60}{m}. Sentuhlah pada \qty{20}{cm}
(atau \qty{40}{cm}) dari sebuah ujung: simpul di sana hanya menyisakan
$n = 3, 6, \dots$, dan senarnya berbunyi pada \qty{600}{Hz}.
\end{solution}

\begin{solution}{exo:b1:wave-propagation:6}
$\lambda = \qty{0.40}{m}$. Pada ruasnya, pada jarak $x$ dari pengeras
suara 1, $\delta = (2.0 - x) - x = 2.0 - 2x$; senyap untuk $\delta = \pm\lambda/2,
\pm 3\lambda/2, \dots$: $x = 0.1, 0.3, \dots, 1.9$ m, setiap \qty{20}{cm};
titik tengahnya keras. Pada sumbu baginya $\delta = 0$: keras di
sepanjangnya (hanya ada peluruhan $1/r$). Berlawanan \omterm{def:b1:wave-propagation:signal}{fase}: yang keras
dan yang senyap bertukar --- senyap di titik tengahnya dan setiap \qty{20}{cm} darinya.
\end{solution}

\begin{solution}{exo:b1:wave-propagation:7}
$i = \lambda D/a = 5.89\times10^{-7} \times 1.5/2.0\times10^{-4} =
\qty{4.4}{mm}$; biru: \qty{3.4}{mm}. Dalam $\pm\qty{1.0}{cm}$: pita di
$0$, $\pm 4.4$, $\pm 8.8$ mm --- lima buah.
\end{solution}

\begin{solution}{exo:b1:wave-propagation:8}
Sama besar: $I_{\max} = 4I_0$, $I_{\min} = 0$, kontrasnya $1$. $I_0$ dan $I_0/4$:
$I_{\max} = I_0(1 + \tfrac14 + 2 \times \tfrac12) = 2.25I_0$, $I_{\min} =
I_0(1.25 - 1) = 0.25I_0$; kontrasnya $2.0/2.5 = 0.8$.
\end{solution}

\begin{solution}{exo:b1:wave-propagation:9}
Tertutup--terbuka: $f = (2n + 1)c/4L = 170$, $510$, \qty{850}{Hz} (hanya
harmonik ganjil: klarinetnya). Terbuka--terbuka: $nc/2L = 340$, $680$,
\qty{1020}{Hz} (serulingnya).
\end{solution}

\begin{solution}{exo:b1:wave-propagation:10}
$v_\varphi = \sqrt{9.81 \times 10/2\pi} = \qty{3.9}{m/s}$ dan
\qty{12.5}{m/s} untuk \qty{100}{m}. Alun panjangnya lebih dulu: $10^6/12.5 =
\qty{22}{h}$ terhadap $10^6/3.95 = \qty{70}{h}$, dua hari lebih awal.
Lajunya bergantung pada panjang gelombang: lautan bercampur menyortir
dirinya menurut panjang gelombang sambil berjalan --- itulah dispersi.
\end{solution}

\begin{solution}{exo:b1:wave-propagation:11}
Moda $n$ berbentuk $\sin(n\pi x/L)$. Simpul yang dipaksakan di $x = L/3$
menuntut $\sin(n\pi/3) = 0$, yaitu $n$ kelipatan $3$: frekuensinya $3f_1,
6f_1, \dots$ Memetik di $L/2$ memberi senarnya simpangan maksimum di
sana, yang hanya dapat disumbang moda $n$ jika $\sin(n\pi/2) \neq 0$:
moda genap, yang bersimpul di $L/2$, tak menerima apa pun.
\end{solution}

\begin{solution}{exo:b1:wave-propagation:12}
$\lambda = c/f = \qty{3.0}{m}$. Beda \omterm{def:b1:wave-propagation:signal}{fasenya} $2\pi(2d)/\lambda + \pi$;
destruktif ketika ia sama dengan $(2m + 1)\pi$: $2d/\lambda = m$,
$d = m\lambda/2 = 0, 1.5, 3.0, \dots$ m; konstruktif ketika $2d/\lambda =
m + \tfrac12$: $d = 0.75, 2.25, \dots$ m. Dinding dan atap terowongan
memantulkan gelombangnya: sebuah pola berdiri dengan bintik mati setiap
setengah panjang gelombang menyusuri jalur mobilnya.
\end{solution}

\begin{solution}{pb:b1:wave-propagation:1}
\textbf{1.} $[F/\mu] = M L T^{-2}/(M L^{-1}) = L^2 T^{-2}$: kuadrat
sebuah laju.

\textbf{2.} Gelombang menuju $-x$ bergantung pada $t + x/c$, yaitu pada
$\omega t + kx$; $k = \omega/c = 2\pi f/c$.

\textbf{3.} $c = \sqrt{0.90/1.0\times10^{-3}} = \qty{30}{m/s}$;
$\lambda = c/f = \qty{0.60}{m}$; $k = 2\pi/\lambda = \qty{10.5}{rad/m}$.

\textbf{4.} $\Delta\varphi = kd = 10.5 \times 0.15 = \qty{1.6}{rad} = \pi/2$.

\textbf{5.} $v_{\max} = A\omega = 2.0\times10^{-3} \times 2\pi \times 50 =
\qty{0.63}{m/s}$, lima puluh kali lebih kecil daripada $c$: titik-titik
senarnya bergerak lambat sementara bentuknya melesat.

\textbf{6.} Ia berjalan menuju $+x$: \omterm{def:b1:wave-propagation:signal}{fasenya} $\omega t - kx$. Pada ujung
tetapnya simpangan totalnya lenyap setiap saat: $s_i(0, t) + s_r(0, t)
= A\cos\omega t + (-A)\cos\omega t = 0$.

\textbf{7.} $\cos(\omega t + kx) - \cos(\omega t - kx) =
-2\sin(\omega t)\sin(kx)$, sehingga $s = -2A\sin(kx)\sin(\omega t)$.

\textbf{8.} Simpul: $\sin kx = 0$, $x = n\lambda/2$; perut $x = (n +
\tfrac12)\lambda/2$; simpul bertetangga berjarak $\lambda/2$.

\textbf{9.} $\sin kL = 0$: $L = n\lambda/2$, yaitu $f_n = nc/2L$.

\textbf{10.} $n = 2L/\lambda = 2.4/0.60 = 4$: empat perut, sebuah
bilangan bulat --- resonansi.

\textbf{11.} $3$ perut: $\lambda = 2L/3 = \qty{0.80}{m}$, $c = \lambda f
= \qty{40}{m/s}$, $F = \mu c^2 = \qty{1.6}{N}$. $6$ perut: $\lambda =
\qty{0.40}{m}$, $c = \qty{20}{m/s}$, $F = \qty{0.40}{N}$.

\textbf{12.} $c = 2Lf/n = 2 \times 1.20 \times 50/3 = \qty{40}{m/s}$,
$\mu = F/c^2 = 1.6/1600 = \qty{1.0e-3}{kg/m} = \qty{1.0}{g/m}$.

\textbf{13.} $\mu = mgn^2/(4L^2f^2)$: $u(\mu)/\mu = \sqrt{(u_m/m)^2 +
(2u_L/L)^2} = \sqrt{0.61^2 + 0.33^2}\,\% = 0.7\%$; massanya yang membatasi.

\textbf{14.} Nada dasarnya. $c = 2Lf_1$: E rendah $2 \times 0.650 \times
82.4 = \qty{107}{m/s}$; E tinggi \qty{428}{m/s}.

\textbf{15.} $F = \mu c^2$: E rendah $6.0\times10^{-3} \times 107^2 =
\qty{69}{N}$; E tinggi $4.0\times10^{-4} \times 428^2 = \qty{73}{N}$ ---
tegangan yang sebanding. Menurunkan $f$ dengan faktor $4$ pada $L$ dan
$\mu$ tetap akan menuntut tegangan $16$ kali lebih kecil: senar yang
lembek dan mendengung. Lilitannya menaikkan $\mu$ sebagai gantinya.

\textbf{16.} $f_n = 2^{n/12}f_1$ dan $f \propto 1/L$: $L_n = 2^{-n/12}L$,
jaraknya dari nat $L - L_n = L(1 - 2^{-n/12})$: fret 1 \qty{36.5}{mm},
fret 5 \qty{163}{mm}, fret 7 \qty{216}{mm}, fret 12 \qty{325}{mm}.

\textbf{17.} Tiap jaraknya $2^{-1/12} = 0.944$ kali jarak sebelumnya.

\textbf{18.} Faktornya $2^{3/12} = 1.19$ (tiga seminada); panjangnya
$650 \times 2^{-3/12} = \qty{547}{mm}$.

\textbf{19.} Menyentuh di $L/2$: hanya moda genap, nadanya $2f_1 =
\qty{165}{Hz}$; di $L/3$: kelipatan $3$, nadanya $3f_1 = \qty{247}{Hz}$.

\textbf{20.} Dekat sebuah ujung, $\sin(n\pi x/L) \approx n\pi x/L$
tumbuh bersama $n$: moda tinggi relatif lebih terbangkitkan --- cerah.
Di atas lubang suaranya (dekat tengahnya) moda genapnya lemah dan nada
dasarnya mendominasi --- lembut.

\textbf{21.} $3f_1 = \qty{330}{Hz}$ melayang \qty{2}{Hz} terhadap $2f_2$:
$f_2 = 166$ atau \qty{164}{Hz}, nisbahnya $1.509$ atau $1.491$ (kuint
bertemperamen sama adalah $2^{7/12} = 1.498$).

\textbf{22.} Harmonik keempat, $4f$. Tiga layangan: $4f = 437$ atau
\qty{443}{Hz}, $f = 109.25$ atau \qty{110.75}{Hz}.

\textbf{23.} Layangannya makin cepat ketika dikencangkan, jadi senarnya
terlalu tinggi: \qty{110.75}{Hz}. Ia harus turun sebesar
$0.75/110.75 = 0.68\%$, yaitu tegangannya sebesar $1.4\%$.

\textbf{24.} $f = (n/2L)\sqrt{F/\mu}$, sehingga $\ln f = \tfrac12\ln F + \text{tetap}$
dan $\dd f/f = \tfrac12\,\dd F/F$.

\textbf{25.} $f_1 = \dfrac{1}{2L}\sqrt{\dfrac{F}{\mu}}$: panjangnya
menetapkan $\lambda_1 = 2L$, tegangan dan massanya menetapkan $c$, fret
menskalakan $L$ dengan $2^{-n/12}$, dan layangan menyingkap
$f - f_{\text{acuan}}$. Di balik semuanya, satu hubungan:
$\lambda f = c$, dengan syarat \omterm{thm:b1:wave-propagation:standing}{gelombang berdiri} $L = n\lambda/2$.
\end{solution}
